%(BEGIN_QUESTION)
% Copyright 2006, Tony R. Kuphaldt, released under the Creative Commons Attribution License (v 1.0)
% This means you may do almost anything with this work of mine, so long as you give me proper credit

Et hydraulikksystem har to sylindre koblet sammen med høytrykksrør. Stempeldiameteren i sylinder \#1 er 3 tommer, mens stempeldiameteren i sylinder \#2 er 4.5 tommer. Hvor stor kraft vil stempelet i sylinder \#1 utøve dersom stempelet i sylinder \#2 trykkes med 200 pund kraft? Hvilket fluidtrykk oppstår i hydraulikkrøret når 200 pund kraft påføres stempelet i sylinder \#2?

\vskip 10pt

Bestem også hvilket av de to stemplene som forflytter seg lengst, og forklar hvorfor.

\vskip 20pt \vbox{\hrule \hbox{\strut \vrule{} {\bf Forslag til sokratisk diskusjon} \vrule} \hrule}

\begin{itemize}
\item{} En nyttig problemløsningsmetode er å skissere et enkelt diagram av systemet som analyseres. Dette er nyttig selv når du allerede har et diagram i oppgaven, fordi en enklere skisse ofte reduserer kompleksiteten og gjør løsningen lettere. Tegn din egen skisse av hvordan oppgitte data henger sammen, og bruk den til å forklare løsningen.
\end{itemize}

\underbar{file i00151}
%(END_QUESTION)





%(BEGIN_ANSWER)

Fluidtrykket blir 39.5 PSI, og stempelkraften i sylinder \#1 blir 88.889 pund.

\vskip 10pt

Stempel \#1 vil forflytte seg lengre enn stempel \#2.

%(END_ANSWER)





%(BEGIN_NOTES)






\vfil \eject

\noindent
{\bf Forhåndsquiz:}

Anta at 800 PSI hydraulikktrykk virker på et rundt stempel med diameter 2 tommer. Beregn kraften som dette fluidtrykket genererer på stempelet.

\begin{itemize}
\item{} 2,513.3 pund
\vskip 5pt
\item{} 254.6 pund
\vskip 5pt
\item{} 5,026.5 pund
\vskip 5pt
\item{} 800 pund
\vskip 5pt
\item{} 10,053.1 pund
\vskip 5pt
\item{} 1600 pund
\end{itemize}


%INDEX% Physics, fluids: pressure, force, and area
%INDEX% Physics, static fluids: Pascal's Principle

%(END_NOTES)
